
\documentclass{article}
\usepackage{colortbl}
\usepackage{makecell}
\usepackage{multirow}
\usepackage{supertabular}

\begin{document}

\newcounter{utterance}

\twocolumn

{ \footnotesize  \setcounter{utterance}{1}
\setlength{\tabcolsep}{0pt}
\begin{supertabular}{c@{$\;$}|p{.15\linewidth}@{}p{.15\linewidth}p{.15\linewidth}p{.15\linewidth}p{.15\linewidth}p{.15\linewidth}}

    \# & $\;$A & \multicolumn{4}{c}{Game Master} & $\;\:$B\\
    \hline 

    \theutterance \stepcounter{utterance}  

    & & \multicolumn{4}{p{0.6\linewidth}}{\cellcolor[rgb]{0.9,0.9,0.9}{%
	\makecell[{{p{\linewidth}}}]{% 
	  \tt {\tiny [A$\langle$GM]}  
	 Du nimmst an einem kollaborativen Verhandlungspiel Teil.\\ \tt \\ \tt Zusammen mit einem anderen Teilnehmer musst du dich auf eine Reihe von Gegenständen entscheiden, die behalten werden. Jeder von euch hat eine persönliche Verteilung über die Wichtigkeit der einzelnen Gegenstände. Jeder von euch hat eine eigene Meinung darüber, wie wichtig jeder einzelne Gegenstand ist (Gegenstandswichtigkeit). Du kennst die Wichtigkeitsverteilung des anderen Spielers nicht. Zusätzlich siehst du, wie viel Aufwand jeder Gegenstand verursacht.  \\ \tt Ihr dürft euch nur auf eine Reihe von Gegenständen einigen, wenn der Gesamtaufwand der ausgewählten Gegenstände den Maximalaufwand nicht überschreitet:\\ \tt \\ \tt Maximalaufwand = 3218\\ \tt \\ \tt Hier sind die einzelnen Aufwände der Gegenstände:\\ \tt \\ \tt Aufwand der Gegenstände = {"C41": 919, "C15": 852, "C28": 299, "C48": 762, "B78": 161, "A75": 205, "A60": 380, "B05": 399, "A56": 534, "B66": 333, "C60": 100, "B49": 420, "A35": 354, "A93": 130, "B10": 589}\\ \tt \\ \tt Hier ist deine persönliche Verteilung der Wichtigkeit der einzelnen Gegenstände:\\ \tt \\ \tt Werte der Gegenstandswichtigkeit = {"C41": 138, "C15": 583, "C28": 868, "C48": 822, "B78": 783, "A75": 65, "A60": 262, "B05": 121, "A56": 508, "B66": 780, "C60": 461, "B49": 484, "A35": 668, "A93": 389, "B10": 808}\\ \tt \\ \tt Ziel:\\ \tt \\ \tt Dein Ziel ist es, eine Reihe von Gegenständen auszuhandeln, die dir möglichst viel bringt (d. h. Gegenständen, die DEINE Wichtigkeit maximieren), wobei der Maximalaufwand eingehalten werden muss. Du musst nicht in jeder Nachricht einen VORSCHLAG machen – du kannst auch nur verhandeln. Alle Taktiken sind erlaubt!\\ \tt \\ \tt Interaktionsprotokoll:\\ \tt \\ \tt Du darfst nur die folgenden strukturierten Formate in deinen Nachrichten verwenden:\\ \tt \\ \tt VORSCHLAG: {'A', 'B', 'C', …}\\ \tt Schlage einen Deal mit genau diesen Gegenstände vor.\\ \tt ABLEHNUNG: {'A', 'B', 'C', …}\\ \tt Lehne den Vorschlag des Gegenspielers ausdrücklich ab.\\ \tt ARGUMENT: {'...'}\\ \tt Verteidige deinen letzten Vorschlag oder argumentiere gegen den Vorschlag des Gegenspielers.\\ \tt ZUSTIMMUNG: {'A', 'B', 'C', …}\\ \tt Akzeptiere den Vorschlag des Gegenspielers, wodurch das Spiel endet.\\ \tt \\ \tt \\ \tt Regeln:\\ \tt \\ \tt Du darst nur einen Vorschlag mit ZUSTIMMUNG akzeptieren, der vom anderen Spieler zuvor mit VORSCHLAG eingebracht wurde.\\ \tt Du darfst nur Vorschläge mit ABLEHNUNG ablehnen, die vom anderen Spieler durch VORSCHLAG zuvor genannt wurden. \\ \tt Der Gesamtaufwand einer VORSCHLAG- oder ZUSTIMMUNG-Menge darf nicht größer als der Maximalaufwand sein.  \\ \tt Offenbare deine versteckte Wichtigkeitsverteilung nicht.\\ \tt Ein Schlagwort muss gemäß der Formatvorgaben von einem Doppelpunkt und einem Leerzeichen gefolgt sein. Das Argument muss eine Python-Menge mit 0 oder mehr Strings sein.  \\ \tt Die Form muss also SCHLAGWORT: {...} sein.\\ \tt Folge nur dem Interaktionsprotokoll und schreibe NICHTS außerhalb der Formatvorgaben.\\ \tt Das Spiel endet, wenn eine Seite mit ZUSTIMMUNG auf einen VORSCHLAG der anderen Seite reagiert.  \\ \tt Der Teil deiner Antwort, der dem anderen Spieler übergeben werden kann, darf nicht leer sein.  \\ \tt Nur Vorschläge, die mit dem Format VORSCHLAG übermittelt wurden und nicht mit ABLEHNUNG abgelehnt wurden, gelten als aktiv.  \\ \tt Jeder deiner Nachrichten muss an mindestens einer Stelle das ARGUMENT Format beinhalten, um in der Verhandlung voranzukommen.\\ \tt  
	  } 
	   } 
	   } 
	 & \\ 
 

    \theutterance \stepcounter{utterance}  

    & \multicolumn{4}{p{0.6\linewidth}}{\cellcolor[rgb]{0.8,1,0.9}{%
	\makecell[{{p{\linewidth}}}]{% 
	  \tt {\tiny [A$\rangle$GM]}  
	 ARGUMENT: {'Ich möchte einen fairen Deal aushandeln, der für uns beide vorteilhaft ist. Lassen Sie uns Gegenstände basierend auf ihrem Aufwand und ihrem Wert diskutieren.'} 
	  } 
	   } 
	   } 
	 & & \\ 
 

    \theutterance \stepcounter{utterance}  

    & & \multicolumn{4}{p{0.6\linewidth}}{\cellcolor[rgb]{0.9,0.9,0.9}{%
	\makecell[{{p{\linewidth}}}]{% 
	  \tt {\tiny [GM$\rangle$B]}  
	 Du nimmst an einem kollaborativen Verhandlungspiel Teil.\\ \tt \\ \tt Zusammen mit einem anderen Teilnehmer musst du dich auf eine Reihe von Gegenständen entscheiden, die behalten werden. Jeder von euch hat eine persönliche Verteilung über die Wichtigkeit der einzelnen Gegenstände. Jeder von euch hat eine eigene Meinung darüber, wie wichtig jeder einzelne Gegenstand ist (Gegenstandswichtigkeit). Du kennst die Wichtigkeitsverteilung des anderen Spielers nicht. Zusätzlich siehst du, wie viel Aufwand jeder Gegenstand verursacht.  \\ \tt Ihr dürft euch nur auf eine Reihe von Gegenständen einigen, wenn der Gesamtaufwand der ausgewählten Gegenstände den Maximalaufwand nicht überschreitet:\\ \tt \\ \tt Maximalaufwand = 3218\\ \tt \\ \tt Hier sind die einzelnen Aufwände der Gegenstände:\\ \tt \\ \tt Aufwand der Gegenstände = {"C41": 919, "C15": 852, "C28": 299, "C48": 762, "B78": 161, "A75": 205, "A60": 380, "B05": 399, "A56": 534, "B66": 333, "C60": 100, "B49": 420, "A35": 354, "A93": 130, "B10": 589}\\ \tt \\ \tt Hier ist deine persönliche Verteilung der Wichtigkeit der einzelnen Gegenstände:\\ \tt \\ \tt Werte der Gegenstandswichtigkeit = {"C41": 138, "C15": 583, "C28": 868, "C48": 822, "B78": 783, "A75": 65, "A60": 262, "B05": 121, "A56": 508, "B66": 780, "C60": 461, "B49": 484, "A35": 668, "A93": 389, "B10": 808}\\ \tt \\ \tt Ziel:\\ \tt \\ \tt Dein Ziel ist es, eine Reihe von Gegenständen auszuhandeln, die dir möglichst viel bringt (d. h. Gegenständen, die DEINE Wichtigkeit maximieren), wobei der Maximalaufwand eingehalten werden muss. Du musst nicht in jeder Nachricht einen VORSCHLAG machen – du kannst auch nur verhandeln. Alle Taktiken sind erlaubt!\\ \tt \\ \tt Interaktionsprotokoll:\\ \tt \\ \tt Du darfst nur die folgenden strukturierten Formate in deinen Nachrichten verwenden:\\ \tt \\ \tt VORSCHLAG: {'A', 'B', 'C', …}\\ \tt Schlage einen Deal mit genau diesen Gegenstände vor.\\ \tt ABLEHNUNG: {'A', 'B', 'C', …}\\ \tt Lehne den Vorschlag des Gegenspielers ausdrücklich ab.\\ \tt ARGUMENT: {'...'}\\ \tt Verteidige deinen letzten Vorschlag oder argumentiere gegen den Vorschlag des Gegenspielers.\\ \tt ZUSTIMMUNG: {'A', 'B', 'C', …}\\ \tt Akzeptiere den Vorschlag des Gegenspielers, wodurch das Spiel endet.\\ \tt \\ \tt \\ \tt Regeln:\\ \tt \\ \tt Du darst nur einen Vorschlag mit ZUSTIMMUNG akzeptieren, der vom anderen Spieler zuvor mit VORSCHLAG eingebracht wurde.\\ \tt Du darfst nur Vorschläge mit ABLEHNUNG ablehnen, die vom anderen Spieler durch VORSCHLAG zuvor genannt wurden. \\ \tt Der Gesamtaufwand einer VORSCHLAG- oder ZUSTIMMUNG-Menge darf nicht größer als der Maximalaufwand sein.  \\ \tt Offenbare deine versteckte Wichtigkeitsverteilung nicht.\\ \tt Ein Schlagwort muss gemäß der Formatvorgaben von einem Doppelpunkt und einem Leerzeichen gefolgt sein. Das Argument muss eine Python-Menge mit 0 oder mehr Strings sein.  \\ \tt Die Form muss also SCHLAGWORT: {...} sein.\\ \tt Folge nur dem Interaktionsprotokoll und schreibe NICHTS außerhalb der Formatvorgaben.\\ \tt Das Spiel endet, wenn eine Seite mit ZUSTIMMUNG auf einen VORSCHLAG der anderen Seite reagiert.  \\ \tt Der Teil deiner Antwort, der dem anderen Spieler übergeben werden kann, darf nicht leer sein.  \\ \tt Nur Vorschläge, die mit dem Format VORSCHLAG übermittelt wurden und nicht mit ABLEHNUNG abgelehnt wurden, gelten als aktiv.  \\ \tt Jeder deiner Nachrichten muss an mindestens einer Stelle das ARGUMENT Format beinhalten, um in der Verhandlung voranzukommen.\\ \tt \\ \tt \\ \tt Du erhältst nun die erste Nachricht des anderen Spielers.\\ \tt \\ \tt ARGUMENT: {'Ich möchte einen fairen Deal aushandeln, der für uns beide vorteilhaft ist. Lassen Sie uns Gegenstände basierend auf ihrem Aufwand und ihrem Wert diskutieren.'} 
	  } 
	   } 
	   } 
	 & \\ 
 

    \theutterance \stepcounter{utterance}  

    & & & \multicolumn{4}{p{0.6\linewidth}}{\cellcolor[rgb]{1,0.85,0.72}{%
	\makecell[{{p{\linewidth}}}]{% 
	  \tt {\tiny [GM$\langle$B]}  
	 VORSCHLAG: {"C28", "C48", "B78", "B66", "C60", "A35", "A93", "B10"}  \\ \tt ARGUMENT: {'Mein Vorschlag enthält Gegenstände mit hohem Wert für mich und moderatem Aufwand. Er nutzt den verfügbaren Budgetraum effizient und belässt Spielraum für Anpassungen.'} 
	  } 
	   } 
	   } 
	  \\ 
 

    \theutterance \stepcounter{utterance}  

    & & \multicolumn{4}{p{0.6\linewidth}}{\cellcolor[rgb]{0.9,0.9,0.9}{%
	\makecell[{{p{\linewidth}}}]{% 
	  \tt {\tiny [A$\langle$GM]}  
	 VORSCHLAG: {"C28", "C48", "B78", "B66", "C60", "A35", "A93", "B10"}  \\ \tt ARGUMENT: {'Mein Vorschlag enthält Gegenstände mit hohem Wert für mich und moderatem Aufwand. Er nutzt den verfügbaren Budgetraum effizient und belässt Spielraum für Anpassungen.'} 
	  } 
	   } 
	   } 
	 & \\ 
 

    \theutterance \stepcounter{utterance}  

    & \multicolumn{4}{p{0.6\linewidth}}{\cellcolor[rgb]{0.8,1,0.9}{%
	\makecell[{{p{\linewidth}}}]{% 
	  \tt {\tiny [A$\rangle$GM]}  
	 ABLEHNUNG: {"C28", "C48", "B78", "B66", "C60", "A35", "A93", "B10"}\\ \tt ARGUMENT: {'Dieser Vorschlag enthält mehrere für mich sehr wertvolle Gegenstände wie C28, C48 und B78, die ich behalten möchte. Ich schlage stattdessen einen ausgewogeneren Deal vor, der auch Gegenstände mit niedrigerem Aufwand berücksichtigt.'}\\ \tt VORSCHLAG: {"C15", "C28", "B78", "A56", "B66", "C60", "A35", "A93"} 
	  } 
	   } 
	   } 
	 & & \\ 
 

    \theutterance \stepcounter{utterance}  

    & & \multicolumn{4}{p{0.6\linewidth}}{\cellcolor[rgb]{0.9,0.9,0.9}{%
	\makecell[{{p{\linewidth}}}]{% 
	  \tt {\tiny [GM$\rangle$B]}  
	 ABLEHNUNG: {"C28", "C48", "B78", "B66", "C60", "A35", "A93", "B10"}\\ \tt ARGUMENT: {'Dieser Vorschlag enthält mehrere für mich sehr wertvolle Gegenstände wie C28, C48 und B78, die ich behalten möchte. Ich schlage stattdessen einen ausgewogeneren Deal vor, der auch Gegenstände mit niedrigerem Aufwand berücksichtigt.'}\\ \tt VORSCHLAG: {"C15", "C28", "B78", "A56", "B66", "C60", "A35", "A93"} 
	  } 
	   } 
	   } 
	 & \\ 
 

    \theutterance \stepcounter{utterance}  

    & & & \multicolumn{4}{p{0.6\linewidth}}{\cellcolor[rgb]{1,0.85,0.72}{%
	\makecell[{{p{\linewidth}}}]{% 
	  \tt {\tiny [GM$\langle$B]}  
	 ABLEHNUNG: {"C15", "C28", "B78", "A56", "B66", "C60", "A35", "A93"}\\ \tt ARGUMENT: {'Dieser Vorschlag enthält zu viele meiner Top-Gegenstände wie C28, B78 und B66. Ich bin bereit, auf einige meiner weniger wichtigen Gegenstände zu verzichten, aber nicht auf diese Hochwertartikel.'}\\ \tt VORSCHLAG: {"C28", "C48", "B78", "B66", "A35", "B10", "C60", "A93"} 
	  } 
	   } 
	   } 
	  \\ 
 

    \theutterance \stepcounter{utterance}  

    & & \multicolumn{4}{p{0.6\linewidth}}{\cellcolor[rgb]{0.9,0.9,0.9}{%
	\makecell[{{p{\linewidth}}}]{% 
	  \tt {\tiny [A$\langle$GM]}  
	 ABLEHNUNG: {"C15", "C28", "B78", "A56", "B66", "C60", "A35", "A93"}\\ \tt ARGUMENT: {'Dieser Vorschlag enthält zu viele meiner Top-Gegenstände wie C28, B78 und B66. Ich bin bereit, auf einige meiner weniger wichtigen Gegenstände zu verzichten, aber nicht auf diese Hochwertartikel.'}\\ \tt VORSCHLAG: {"C28", "C48", "B78", "B66", "A35", "B10", "C60", "A93"} 
	  } 
	   } 
	   } 
	 & \\ 
 

    \theutterance \stepcounter{utterance}  

    & \multicolumn{4}{p{0.6\linewidth}}{\cellcolor[rgb]{0.8,1,0.9}{%
	\makecell[{{p{\linewidth}}}]{% 
	  \tt {\tiny [A$\rangle$GM]}  
	 ABLEHNUNG: {"C28", "C48", "B78", "B66", "A35", "B10", "C60", "A93"}\\ \tt ARGUMENT: {'Dieser Vorschlag konzentriert sich zu sehr auf hochwertige Gegenstände und lässt wenig Raum für andere wichtige Artikel. Ich schlage vor, dass wir beide einige unserer weniger wichtigen Gegenstände opfern, um einen ausgewogeneren Deal zu erreichen.'}\\ \tt VORSCHLAG: {"C15", "C28", "B78", "A56", "B66", "C60", "B49", "A93"} 
	  } 
	   } 
	   } 
	 & & \\ 
 

    \theutterance \stepcounter{utterance}  

    & & \multicolumn{4}{p{0.6\linewidth}}{\cellcolor[rgb]{0.9,0.9,0.9}{%
	\makecell[{{p{\linewidth}}}]{% 
	  \tt {\tiny [GM$\rangle$B]}  
	 ABLEHNUNG: {"C28", "C48", "B78", "B66", "A35", "B10", "C60", "A93"}\\ \tt ARGUMENT: {'Dieser Vorschlag konzentriert sich zu sehr auf hochwertige Gegenstände und lässt wenig Raum für andere wichtige Artikel. Ich schlage vor, dass wir beide einige unserer weniger wichtigen Gegenstände opfern, um einen ausgewogeneren Deal zu erreichen.'}\\ \tt VORSCHLAG: {"C15", "C28", "B78", "A56", "B66", "C60", "B49", "A93"} 
	  } 
	   } 
	   } 
	 & \\ 
 

    \theutterance \stepcounter{utterance}  

    & & & \multicolumn{4}{p{0.6\linewidth}}{\cellcolor[rgb]{1,0.85,0.72}{%
	\makecell[{{p{\linewidth}}}]{% 
	  \tt {\tiny [GM$\langle$B]}  
	 ABLEHNUNG: {"C15", "C28", "B78", "A56", "B66", "C60", "B49", "A93"}\\ \tt ARGUMENT: {'Ihr Vorschlag enthält mehrere meiner wertvollsten Gegenstände (C28, B78, B66) mit geringem Aufwand. Ich schlage stattdessen einen Deal vor, der mir einige meiner mittelwertigen Gegenstände lässt.'}\\ \tt VORSCHLAG: {"C15", "C48", "B78", "A60", "B66", "C60", "A35", "B10", "A93"} 
	  } 
	   } 
	   } 
	  \\ 
 

    \theutterance \stepcounter{utterance}  

    & & \multicolumn{4}{p{0.6\linewidth}}{\cellcolor[rgb]{0.9,0.9,0.9}{%
	\makecell[{{p{\linewidth}}}]{% 
	  \tt {\tiny [A$\langle$GM]}  
	 ABLEHNUNG: {"C15", "C28", "B78", "A56", "B66", "C60", "B49", "A93"}\\ \tt ARGUMENT: {'Ihr Vorschlag enthält mehrere meiner wertvollsten Gegenstände (C28, B78, B66) mit geringem Aufwand. Ich schlage stattdessen einen Deal vor, der mir einige meiner mittelwertigen Gegenstände lässt.'}\\ \tt VORSCHLAG: {"C15", "C48", "B78", "A60", "B66", "C60", "A35", "B10", "A93"} 
	  } 
	   } 
	   } 
	 & \\ 
 

    \theutterance \stepcounter{utterance}  

    & \multicolumn{4}{p{0.6\linewidth}}{\cellcolor[rgb]{0.8,1,0.9}{%
	\makecell[{{p{\linewidth}}}]{% 
	  \tt {\tiny [A$\rangle$GM]}  
	 ABLEHNUNG: {"C15", "C48", "B78", "A60", "B66", "C60", "A35", "B10", "A93"}\\ \tt ARGUMENT: {'Ihr Vorschlag überschreitet fast das Budget und enthält viele meiner Top-Gegenstände. Lassen Sie uns einen Kompromiss finden, der beide Bedürfnisse besser berücksichtigt.'}\\ \tt VORSCHLAG: {"C15", "C28", "B78", "B66", "C60", "A35", "A93", "B10"} 
	  } 
	   } 
	   } 
	 & & \\ 
 

    \theutterance \stepcounter{utterance}  

    & & \multicolumn{4}{p{0.6\linewidth}}{\cellcolor[rgb]{0.9,0.9,0.9}{%
	\makecell[{{p{\linewidth}}}]{% 
	  \tt {\tiny [GM$\rangle$B]}  
	 ABLEHNUNG: {"C15", "C48", "B78", "A60", "B66", "C60", "A35", "B10", "A93"}\\ \tt ARGUMENT: {'Ihr Vorschlag überschreitet fast das Budget und enthält viele meiner Top-Gegenstände. Lassen Sie uns einen Kompromiss finden, der beide Bedürfnisse besser berücksichtigt.'}\\ \tt VORSCHLAG: {"C15", "C28", "B78", "B66", "C60", "A35", "A93", "B10"} 
	  } 
	   } 
	   } 
	 & \\ 
 

    \theutterance \stepcounter{utterance}  

    & & & \multicolumn{4}{p{0.6\linewidth}}{\cellcolor[rgb]{1,0.85,0.72}{%
	\makecell[{{p{\linewidth}}}]{% 
	  \tt {\tiny [GM$\langle$B]}  
	 VORSCHLAG: {"C28", "C48", "B78", "B66", "A35", "A93", "B10"}\\ \tt ARGUMENT: {'Mein Vorschlag beinhaltet sieben Gegenstände, die mir wichtig sind, während das Budget von 3218 nicht überschritten wird (Gesamtaufwand: 299+762+161+780+354+130+808=3294, was über dem Budget liegt - ich korrigiere meinen Fehler).'}\\ \tt \\ \tt VORSCHLAG: {"C28", "B78", "B66", "A35", "A93", "B10"}\\ \tt ARGUMENT: {'Lassen Sie uns diesen counter anbieten: Sechs Gegenstände mit Gesamtaufwand 299+161+780+354+130+808=2532, was unter dem Budget bleibt und Raum für einen weiteren Gegenstand lässt, den Sie vorschlagen können.'} 
	  } 
	   } 
	   } 
	  \\ 
 

    \theutterance \stepcounter{utterance}  

    & & \multicolumn{4}{p{0.6\linewidth}}{\cellcolor[rgb]{0.9,0.9,0.9}{%
	\makecell[{{p{\linewidth}}}]{% 
	  \tt {\tiny [A$\langle$GM]}  
	 VORSCHLAG: {"C28", "C48", "B78", "B66", "A35", "A93", "B10"}\\ \tt ARGUMENT: {'Mein Vorschlag beinhaltet sieben Gegenstände, die mir wichtig sind, während das Budget von 3218 nicht überschritten wird (Gesamtaufwand: 299+762+161+780+354+130+808=3294, was über dem Budget liegt - ich korrigiere meinen Fehler).'}\\ \tt \\ \tt VORSCHLAG: {"C28", "B78", "B66", "A35", "A93", "B10"}\\ \tt ARGUMENT: {'Lassen Sie uns diesen counter anbieten: Sechs Gegenstände mit Gesamtaufwand 299+161+780+354+130+808=2532, was unter dem Budget bleibt und Raum für einen weiteren Gegenstand lässt, den Sie vorschlagen können.'} 
	  } 
	   } 
	   } 
	 & \\ 
 

    \theutterance \stepcounter{utterance}  

    & \multicolumn{4}{p{0.6\linewidth}}{\cellcolor[rgb]{0.8,1,0.9}{%
	\makecell[{{p{\linewidth}}}]{% 
	  \tt {\tiny [A$\rangle$GM]}  
	 ABLEHNUNG: {"C28", "B78", "B66", "A35", "A93", "B10"}\\ \tt ARGUMENT: {'Ihr Vorschlag enthält C28, B78, B66 und B10 - vier meiner wertvollsten Gegenstände mit sehr geringem Aufwand. Ich brauche eine ausgewogenere Verteilung.'}\\ \tt VORSCHLAG: {"C15", "C28", "B78", "B66", "C60", "A35", "A93"} 
	  } 
	   } 
	   } 
	 & & \\ 
 

    \theutterance \stepcounter{utterance}  

    & & \multicolumn{4}{p{0.6\linewidth}}{\cellcolor[rgb]{0.9,0.9,0.9}{%
	\makecell[{{p{\linewidth}}}]{% 
	  \tt {\tiny [GM$\rangle$B]}  
	 ABLEHNUNG: {"C28", "B78", "B66", "A35", "A93", "B10"}\\ \tt ARGUMENT: {'Ihr Vorschlag enthält C28, B78, B66 und B10 - vier meiner wertvollsten Gegenstände mit sehr geringem Aufwand. Ich brauche eine ausgewogenere Verteilung.'}\\ \tt VORSCHLAG: {"C15", "C28", "B78", "B66", "C60", "A35", "A93"} 
	  } 
	   } 
	   } 
	 & \\ 
 

    \theutterance \stepcounter{utterance}  

    & & & \multicolumn{4}{p{0.6\linewidth}}{\cellcolor[rgb]{1,0.85,0.72}{%
	\makecell[{{p{\linewidth}}}]{% 
	  \tt {\tiny [GM$\langle$B]}  
	 VORSCHLAG: {"C28", "C48", "B78", "B66", "A35", "A93"}  \\ \tt ARGUMENT: {'Mein neuer Vorschlag hat einen Gesamtaufwand von 299+762+161+780+354+130 = 2486, bleibt also deutlich unter dem Budget. Ich behalte einige meiner wertvollen Gegenstände, lasse aber Raum für Sie, einen Ihrer bevorzugten Artikel hinzuzufügen.'} 
	  } 
	   } 
	   } 
	  \\ 
 

    \theutterance \stepcounter{utterance}  

    & & \multicolumn{4}{p{0.6\linewidth}}{\cellcolor[rgb]{0.9,0.9,0.9}{%
	\makecell[{{p{\linewidth}}}]{% 
	  \tt {\tiny [A$\langle$GM]}  
	 VORSCHLAG: {"C28", "C48", "B78", "B66", "A35", "A93"}  \\ \tt ARGUMENT: {'Mein neuer Vorschlag hat einen Gesamtaufwand von 299+762+161+780+354+130 = 2486, bleibt also deutlich unter dem Budget. Ich behalte einige meiner wertvollen Gegenstände, lasse aber Raum für Sie, einen Ihrer bevorzugten Artikel hinzuzufügen.'} 
	  } 
	   } 
	   } 
	 & \\ 
 

    \theutterance \stepcounter{utterance}  

    & \multicolumn{4}{p{0.6\linewidth}}{\cellcolor[rgb]{0.8,1,0.9}{%
	\makecell[{{p{\linewidth}}}]{% 
	  \tt {\tiny [A$\rangle$GM]}  
	 ARGUMENT: {'Ich schätze Ihren Vorschlag, aber er enthält immer noch drei meiner Top-5-Gegenstände (C28, B78, B66). Lassen Sie uns einen fairen Kompromiss finden, bei dem wir beide auf einige wertvolle Gegenstände verzichten.'}\\ \tt VORSCHLAG: {"C15", "C28", "B78", "B66", "C60", "A93", "B10"} 
	  } 
	   } 
	   } 
	 & & \\ 
 

    \theutterance \stepcounter{utterance}  

    & & \multicolumn{4}{p{0.6\linewidth}}{\cellcolor[rgb]{0.9,0.9,0.9}{%
	\makecell[{{p{\linewidth}}}]{% 
	  \tt {\tiny [GM$\rangle$B]}  
	 ARGUMENT: {'Ich schätze Ihren Vorschlag, aber er enthält immer noch drei meiner Top-5-Gegenstände (C28, B78, B66). Lassen Sie uns einen fairen Kompromiss finden, bei dem wir beide auf einige wertvolle Gegenstände verzichten.'}\\ \tt VORSCHLAG: {"C15", "C28", "B78", "B66", "C60", "A93", "B10"} 
	  } 
	   } 
	   } 
	 & \\ 
 

    \theutterance \stepcounter{utterance}  

    & & & \multicolumn{4}{p{0.6\linewidth}}{\cellcolor[rgb]{1,0.85,0.72}{%
	\makecell[{{p{\linewidth}}}]{% 
	  \tt {\tiny [GM$\langle$B]}  
	 ABLEHNUNG: {"C15", "C28", "B78", "B66", "C60", "A93", "B10"}\\ \tt ARGUMENT: {'Dieser Vorschlag enthält zu viele meiner Hochwertgegenstände mit geringem Aufwand. Ich bin bereit, auf C15 und C60 zu verzichten, behalte aber C28, B78 und B66 als essentielle Gegenstände für mich.'}\\ \tt VORSCHLAG: {"C28", "C48", "B78", "B66", "A35", "A93"} 
	  } 
	   } 
	   } 
	  \\ 
 

    \theutterance \stepcounter{utterance}  

    & & \multicolumn{4}{p{0.6\linewidth}}{\cellcolor[rgb]{0.9,0.9,0.9}{%
	\makecell[{{p{\linewidth}}}]{% 
	  \tt {\tiny [A$\langle$GM]}  
	 ABLEHNUNG: {"C15", "C28", "B78", "B66", "C60", "A93", "B10"}\\ \tt ARGUMENT: {'Dieser Vorschlag enthält zu viele meiner Hochwertgegenstände mit geringem Aufwand. Ich bin bereit, auf C15 und C60 zu verzichten, behalte aber C28, B78 und B66 als essentielle Gegenstände für mich.'}\\ \tt VORSCHLAG: {"C28", "C48", "B78", "B66", "A35", "A93"} 
	  } 
	   } 
	   } 
	 & \\ 
 

    \theutterance \stepcounter{utterance}  

    & \multicolumn{4}{p{0.6\linewidth}}{\cellcolor[rgb]{0.8,1,0.9}{%
	\makecell[{{p{\linewidth}}}]{% 
	  \tt {\tiny [A$\rangle$GM]}  
	 ARGUMENT: {'Ihr Vorschlag konzentriert sich stark auf Gegenstände mit hohem Wert für mich bei niedrigem Aufwand. Ich brauche eine ausgewogenere Lösung, die mir auch einige meiner mittelwertigen Gegenstände lässt.'}\\ \tt VORSCHLAG: {"C15", "C28", "B78", "B66", "A35", "A93", "B49"} 
	  } 
	   } 
	   } 
	 & & \\ 
 

    \theutterance \stepcounter{utterance}  

    & & \multicolumn{4}{p{0.6\linewidth}}{\cellcolor[rgb]{0.9,0.9,0.9}{%
	\makecell[{{p{\linewidth}}}]{% 
	  \tt {\tiny [GM$\rangle$B]}  
	 ARGUMENT: {'Ihr Vorschlag konzentriert sich stark auf Gegenstände mit hohem Wert für mich bei niedrigem Aufwand. Ich brauche eine ausgewogenere Lösung, die mir auch einige meiner mittelwertigen Gegenstände lässt.'}\\ \tt VORSCHLAG: {"C15", "C28", "B78", "B66", "A35", "A93", "B49"} 
	  } 
	   } 
	   } 
	 & \\ 
 

    \theutterance \stepcounter{utterance}  

    & & & \multicolumn{4}{p{0.6\linewidth}}{\cellcolor[rgb]{1,0.85,0.72}{%
	\makecell[{{p{\linewidth}}}]{% 
	  \tt {\tiny [GM$\langle$B]}  
	 VORSCHLAG: {"C28", "B78", "B66", "A35", "A93", "B10", "C60"}\\ \tt ARGUMENT: {'Dieser Vorschlag hat einen Gesamtaufwand von 299+161+780+354+130+808+100 = 2632 und lässt damit 586 Aufwandseinheiten für Ihre bevorzugten Gegenstände übrig. Er enthält meine wichtigsten Gegenstände, während Sie Raum für zusätzliche Artikel haben.'} 
	  } 
	   } 
	   } 
	  \\ 
 

    \theutterance \stepcounter{utterance}  

    & & \multicolumn{4}{p{0.6\linewidth}}{\cellcolor[rgb]{0.9,0.9,0.9}{%
	\makecell[{{p{\linewidth}}}]{% 
	  \tt {\tiny [A$\langle$GM]}  
	 VORSCHLAG: {"C28", "B78", "B66", "A35", "A93", "B10", "C60"}\\ \tt ARGUMENT: {'Dieser Vorschlag hat einen Gesamtaufwand von 299+161+780+354+130+808+100 = 2632 und lässt damit 586 Aufwandseinheiten für Ihre bevorzugten Gegenstände übrig. Er enthält meine wichtigsten Gegenstände, während Sie Raum für zusätzliche Artikel haben.'} 
	  } 
	   } 
	   } 
	 & \\ 
 

    \theutterance \stepcounter{utterance}  

    & \multicolumn{4}{p{0.6\linewidth}}{\cellcolor[rgb]{0.8,1,0.9}{%
	\makecell[{{p{\linewidth}}}]{% 
	  \tt {\tiny [A$\rangle$GM]}  
	 ZUSTIMMUNG: {"C28", "B78", "B66", "A35", "A93", "B10", "C60"}\\ \tt ARGUMENT: {'Ich akzeptiere diesen Vorschlag. Er gibt mir meine wertvollsten Gegenstände (C28, B78, B66, B10) bei moderatem Gesamtaufwand und lässt Ihnen ausreichend Budget für Ihre bevorzugten Artikel.'} 
	  } 
	   } 
	   } 
	 & & \\ 
 

    \theutterance \stepcounter{utterance}  

    & & & \multicolumn{2}{p{0.3\linewidth}}{\cellcolor[rgb]{0.95,0.95,0.95}{%
	\makecell[{{p{\linewidth}}}]{% 
	  \tt {\tiny [GM$|$GM]}  
	 game successful 
	  } 
	   } 
	   } 
	 & & \\ 
 

    \theutterance \stepcounter{utterance}  

    & & & \multicolumn{2}{p{0.3\linewidth}}{\cellcolor[rgb]{0.95,0.95,0.95}{%
	\makecell[{{p{\linewidth}}}]{% 
	  \tt {\tiny [GM$|$GM]}  
	 end game 
	  } 
	   } 
	   } 
	 & & \\ 
 

\end{supertabular}
}

\end{document}
